% questions.tex —— 1977 年·江苏淮阴地区初试卷 · 题目内容片段（无答案）
% 由 tools/generate_projects.py 依来源总账自 1977/1977_数学真题_江苏淮阴地区初试_试题.md 生成，勿手改（改源 MD 后重跑）


\gkpart{江苏淮阴地区初试卷}

\begin{gktip}
\textbf{来源说明}：本题面依据固定源 PDF《1977江苏淮阴地区初试》第 1—2 页逐题转录；第 1 页原书页码 243，第 2 页无页码。源页为试题汇编材料中的淮阴地区初试（预选）卷，仅题面、无解答；页首标题「八、淮阴地区」系汇编章节标题，按原卷逐字保留为本文 H1。原卷大题以中文数字「一、」至「九、」编号，本库转写为 \textbf{1．}至 \textbf{9．}；第 1、4、7 题下的小题原卷以「（1）（2）」编号，照录为（1）（2）。第 1(2) 题原卷指数 $3^{0\cdot 5}$ 以居中圆点作小数分隔（旧排印习惯），本文转写为 $3^{0.5}$。原卷无附图（第 5 题「需画图」系要求考生作图，非原卷附图），无附加题/参考题；配套答案与解析为本库编者独立推导，\textbf{非官方标准答案}。
\end{gktip}

\par\medskip

\begin{question}
\gkti{1}计算：

（1） $4(-1)^{5}\times0.25+8(-0.5^{3})-0.4\div(-0.2^{2})\text{\text{；}}$

（2） $\left[3^{\frac{1}{2}}+\left(\dfrac{1}{2^{-\frac{2}{8}}}\right)^{\frac{2}{3}}\right]\cdot\left[3^{0.5}-(2\sqrt{2})^{\frac{2}{8}}\right]\text{。}$

\end{question}
\begin{question}
\gkti{2}化简：

\nopagebreak[4]
\[
\frac{(a^{2}+1)^{2}-4a^{2}}{-2a^{4}+4a^{2}-2}.
\]

\end{question}
\begin{question}
\gkti{3}求 $\cos\left(-\dfrac{16\pi}{3}\right)\cdot\sin\dfrac{16\pi}{3}$ 的值。

\end{question}
\begin{question}
\gkti{4}

（1） 把 $P$ 点的直角坐标 $\left(\sqrt{3},\ -1\right)$ 化成极坐标；

（2） 把极坐标方程 $\rho=2\sin\theta+3\cos\theta$ 化成直角坐标方程。

\end{question}
\begin{question}
\gkti{5}以直角三角形 $ABC$ 的直角边 $BC$ 为直径作圆交斜边 $AB$ 于 $D$，$E$ 是直角边 $AC$ 的中点，连结 $DE$．则 $DE$ 是此圆的切线。（需画图）

\end{question}
\begin{question}
\gkti{6}用长为 8 米的木料做一个“目”字形窗框，问窗框的宽和高各多少才能使通过的光线最多？

\end{question}
\begin{question}
\gkti{7}计算

（1） $\sqrt[3]{10^{3-\frac{1}{2}\lg 64}}$；

（2） $\log_{2-\sqrt{3}}\left(2+\sqrt{3}\right)$

\end{question}
\begin{question}
\gkti{8}有两点 $A(-2,\ 2)$ 和 $B(2,\ -1)$，试在 $y$ 轴上找出点 $P$，使 $\angle APB=90^{\circ}$。

\end{question}
\begin{question}
\gkti{9}设在椭圆 $\dfrac{x^{2}}{a^{2}}+\dfrac{y^{2}}{b^{2}}=1$ 上任意一点 $P_{1}\left(x_{1},\ y_{1}\right)$ 的切线与椭圆长轴的两端点 $A'$ 和 $A$ 的切线相交于 $Q'$ 和 $Q$，求证：$A'Q'\cdot AQ=b^{2}$。

\end{question}